\begin{frame}{라그랑주 승수법 $\to$ 고유값 문제}
  제약 $v^{T}v = 1$을 승수 $\lambda$로 목적함수에 붙인다:
  \[
    \mathcal{L}(v, \lambda) = v^{T}\Sigma v - \lambda\,(v^{T}v - 1)
  \]
  $v$로 미분하여 0으로 놓는다
  ($\frac{\partial}{\partial v}(v^{T}\Sigma v) = 2\Sigma v$,
  $\Sigma$는 대칭이므로):
  \[
    \frac{\partial \mathcal{L}}{\partial v}
    = 2\Sigma v - 2\lambda v = 0
    \quad\Longrightarrow\quad
    \Sigma v = \lambda v
  \]
  \vskip0.6em
  \begin{block}{이것이 바로 고유값 방정식}
    분산을 최대화하는 방향 $v$는 공분산 행렬의 \textbf{고유벡터}여야
    하고, 그 분산값은 대응하는 \textbf{고유값} $\lambda$다.
    \par 최적점을 원래 목적함수에 대입하면
    $v^{T}\Sigma v = v^{T}(\lambda v) = \lambda(v^{T}v) = \lambda$ ---
    ``그 방향의 분산 = 고유값''을 직접 확인할 수 있다.
  \end{block}
\end{frame}
